% from ega0-18-fr.tex
\oldpage[0\textsubscript{IV-1}]{52}
dans le $A$-anneau $(C,\sigma)$, c'est aussi un $A'$-homomorphisme du
$A'$-anneau $(B,\rho\circ f)$ dans le $A'$-anneau $(C,\sigma\circ f)$~; on
définit ainsi un foncteur canonique de la catégorie des $A$-anneaux dans celle
des $A'$-anneaux.
\end{env}

\begin{env}[18.1.2]
\label{0.18.1.2-fr}
Soient $B$, $E$, $F$ trois $A$-anneaux, $f:E\to B$, $g:F\to B$ deux
$A$-homomorphismes. Rappelons que l'on appelle \emph{produit fibré} de $E$ et
$F$ sur $B$ (pour les $A$-homomorphismes $f$ et $g$) et que l'on note
$E\times_BF$ le sous-anneau $G$ de l'anneau produit $E\times F$, formé des
couples $(x,y)$ tels que $f(x)=g(y)$~; les restrictions $p_1:G\to E$,
$p_2:G\to F$ des projections $\operatorname{pr}_1$ et
$\operatorname{pr}_2$ à $G$ sont encore appelées les \emph{projections
canoniques}~; la structure de $A$-anneau de $G$ est définie par
l'homomorphisme $\alpha\mapsto(\rho(\alpha),\sigma(\alpha))$ (où
$\rho:A\to E$ et $\sigma:A\to F$ sont les homomorphismes structuraux), qui
applique effectivement $A$ dans $G$ en vertu de (18.1.1). La propriété
caractéristique de $E\times_BF$ est que, pour tout couple de
$A$-homomorphismes $u:C\to E$, $v:C\to F$ tels que
$f\circ u=g\circ v$, il existe un $A$-homomorphisme et un seul $w:C\to G$
tel que $u=p_1\circ w$, $v=p_2\circ w$. On peut encore dire que
$E\times_BF$ est \emph{limite projective} du système projectif formé de $B$,
$E$, $F$ et des $A$-homomorphismes $f$, $g$, dans la catégorie des
$A$-anneaux ($0_{\mathrm{III}}$, 8.1.9).
\end{env}

\begin{env}[18.1.3]
\label{0.18.1.3-fr}
Soit $\mathfrak{J}$ l'idéal bilatère de $E$, noyau de $f$~; il est immédiat
que le noyau $\mathfrak{J}'$ de $p_2:G\to F$ est l'idéal formé des éléments
$(x,0)$ où $x\in\mathfrak{J}$~; la restriction
$i_1:\mathfrak{J}'\to\mathfrak{J}$ de $p_1$ est donc un
\emph{isomorphisme} d'anneau \emph{sans élément unité}, et aussi un
isomorphisme de $A$-bimodules. De même, si $\mathfrak{R}$ est le noyau de
$g$, le noyau $\mathfrak{R}'$ de $p_1:G\to E$ est l'idéal des éléments
$(0,y)$ où $y\in\mathfrak{R}$, et la restriction
$i_2:\mathfrak{R}'\to\mathfrak{R}$ de $p_2$ est un isomorphisme (au même
sens). Enfin, il est clair que le noyau du $A$-homomorphisme
$q=f\circ p_1=g\circ p_2$ de $E\times_BF$ dans $B$ est
$\mathfrak{J}\times\mathfrak{R}=\mathfrak{J}'\oplus\mathfrak{R}'$, de sorte
que l'on a le diagramme commutatif
\begin{equation}
\label{0.18.1.3.1-fr}
\tag{18.1.3.1}
  \xymatrix@C=2.2pc@R=1.8pc{
    0\ar[dr] & & 0\ar[d] & 0\ar[d]\\
    & \mathfrak{J}'\oplus\mathfrak{R}'\ar[r]\ar[d]\ar[dr]
      & \mathfrak{R}'\ar[r]^{i_2}_{\sim}\ar[d]^j
      & \mathfrak{R}\ar[d]^{p_2}\\
    0\ar[r] & \mathfrak{J}'\ar[r]\ar[d]^{i_1}_{\sim}
      & G\ar[r]\ar[d]^{p_1}\ar[dr]^q & F\ar[d]^g\\
    0\ar[r] & \mathfrak{J}\ar[r] & E\ar[r]_f & B
  }
\end{equation}

Les définitions et résultats de (18.1.2) et (18.1.3) s'étendent aussitôt au
produit fibré d'une \emph{famille} quelconque $(E_\lambda)_{\lambda\in I}$ de
$A$-anneaux défini par une famille de $A$-homomorphismes
$f_\lambda:E_\lambda\to B$. Nous laissons la formulation de ces résultats au
lecteur.
\end{env}

\begin{env}[18.1.4]
\label{0.18.1.4-fr}
Conformément à la terminologie de (M, VIII), nous appellerons
\emph{$A$-anneau augmenté sur $B$} un $A$-anneau $E$ muni d'un
$A$-homomorphisme \emph{surjectif} (dit \emph{augmentation} de $E$),
$f:E\to B$~; le noyau $\mathfrak{J}$ de $f$ est appelé
\emph{l'idéal d'augmentation}. On dit que le $A$-anneau augmenté $E$ est
\emph{trivial} s'il existe un $A$-homomorphisme de $A$-anneaux $s:B\to E$
qui est \emph{inverse à droite} de l'augmentation $f:E\to B$ (autrement dit
$f\circ s=1_B$). La suite exacte de $A$-bimodules
\[
  0\longrightarrow\mathfrak{J}\longrightarrow E
  \xrightarrow{f}B\longrightarrow0
\]
